% TOP_PROBLEM: {"id": 6200003, "title": "Boundaries of Groups and Kleinian Groups — Problem 3", "category_id": 6, "category": {"id": 6, "name": "geometry", "display_name": "Geometry", "description": "Euclidean and non-Euclidean geometry, geometric structures.", "slug": "geometry", "order_index": 6, "created_at": "2026-07-31T15:26:25.671Z"}, "status": "partially_solved", "problem_number": "AMR-061-0003", "set_id": 13}
% FIRST_POSTED: 2026-10-02
% ENTRY_KIND: statement_only
% UnsolvedMath ID 6200003; source catalogue code AMR-061-0003
% !TeX program = lualatex
% Definition and sources only; this file is not an LLM solution attempt.
\documentclass[11pt,a4paper]{article}
\usepackage[margin=25mm]{geometry}
\usepackage{fontspec}
\setmainfont{lmroman10-regular.otf}[BoldFont=lmroman10-bold.otf,
 ItalicFont=lmroman10-italic.otf,BoldItalicFont=lmroman10-bolditalic.otf]
\usepackage{amsmath,amssymb,mathtools}
\usepackage{unicode-math}
\setmathfont{latinmodern-math.otf}
\AtBeginDocument{\let\setminus\smallsetminus}
\usepackage{xurl,hyperref}
\hypersetup{colorlinks=true,urlcolor=blue}
\setcounter{secnumdepth}{0}
\setlength{\parindent}{0pt}
\setlength{\parskip}{5pt}
\setlength{\emergencystretch}{3em}
\begin{document}
\section*{Boundaries of Groups and Kleinian Groups — Problem 3}\label{6200003}
{\small UnsolvedMath ID 6200003 · AMR-061-0003 · Geometry · AMR Open Problem Lists}
\subsection{Definitions and mathematical statement}
A finitely generated group $G$ is word-hyperbolic if a Cayley graph has $\delta$-thin geodesic triangles for some $\delta\geq0$. Its boundary $\partial G$ consists of geodesic rays starting at a fixed vertex, with two rays identified when their Hausdorff distance is finite. The topology records increasingly long fellow-travelling initial segments. Its homeomorphism type is independent of the finite generating set.

Classify the compact metrizable spaces $Z$ of covering dimension two which occur as
\[
 Z\cong\partial G
\]
for word-hyperbolic groups $G$. Covering dimension two means that every finite open cover has an open refinement of multiplicity at most three, and that no bound of two works for every cover.

The source suggests first imposing additional restrictions, such as no local cut points, no separating arcs or separating Cantor subsets, or no nontrivial splitting in any finite-index subgroup. A cut point of a connected space is a point whose deletion disconnects it; local cut points have this property inside a connected neighborhood. An arc is a subspace homeomorphic to $[0,1]$, and a Cantor set is a nonempty perfect totally disconnected compact metric space.

These restrictions are proposed subcases, not prerequisites for the full classification. The objective concerns homeomorphism types, so an abstract construction of many groups without identifying their boundary spaces does not itself give the requested topological description.
\subsection{Short English statement}
Which two-dimensional compact spaces can occur as boundaries of hyperbolic groups?
\subsection{Sources}
\begin{itemize}
\item[S1] ULAM.AI and the UnsolvedMath Contributors, supplied problems.json, record 6200003 (AMR-061-0003), AMR Open Problem Lists. Catalogue identity and source transcription; CC BY 4.0.
\url{https://huggingface.co/datasets/ulamai/UnsolvedMath}.
\item[S2] Kapovich - Problems on Boundaries of Groups and Kleinian Groups (2005), Problem 3, PDF page 3. Original source indexed by AMR-061-0003.
\url{https://www.math.ucdavis.edu/~kapovich/EPR/problems.pdf}.
\end{itemize}
\end{document}
